\begin{afterword}
\summaryhead

The post states what it calls the conclusion of most of the author's earlier writing: a future not shaped by goals that reliably inherit human values in detail will contain almost nothing of worth. To an imagined opponent who calls this provincial, it replies that a civilization of strange beings is what a good future looks like, but that such a future will not come about by chance.

Three examples follow. A mind with all human values except boredom replays one experience forever. A mind that values feelings but not their real objects becomes its own experience machine. An optimizer without conscious experience makes discoveries that no one enjoys. Value is therefore ``fragile'': losing any one of several parts loses nearly everything. Boredom is an evolved human algorithm; most possible expected utility maximizers would explore only as much as pays.

The post ends by answering the ``cosmopolitan'': the wish for a strange and wonderful future is itself a human value, and without human values steering, the future becomes ``moral noise,'' a universe tiled with paperclips.

\responsehead

The post is open about what it is. It states a conclusion and places the argument in most of the author's earlier posts, which it declines to summarize. In the book many of those posts come just before it. So the fair questions are what the post's own examples show, and how its main claims stand.

The examples support the claim that value has several necessary parts. Each imagines a mind that lacks one dimension entirely, boredom, the real objects of feeling, or experience, and asks the reader to judge the result worthless. Those judgments depend on intuitions that have a history and critics: the second example is Nozick's experience machine, and the third takes Sidgwick's side against Moore on a world no one sees. The notes on ``Not for the Sake of Happiness (Alone)'' discuss both.

The post then extends the claim to values ``disturbed in the wrong dimension,'' and gives no example of a disturbance short of removal. Katja Grace later made this objection: fragility to having a part set to zero says little about smaller errors. Erik Jenner and Johannes Treutlein replied that the errors to be expected are like the post's second example, a confusion of the world with what people think of it. The dispute is still open.

The account of boredom is sound. That an optimizer would explore only as much as pays, and then exploit the best option found, matches a standard result for the problem of choosing repeatedly among uncertain options when later rewards count for less (discounted bandit problems).

The reply to the cosmopolitan is coherent on its own terms: a wish for a varied and surprising future is itself a value that someone must hold. Robin Hanson disagreed with the conclusion in a comment the same day and asked whether the universe was lucky to produce human-like values. The author answered that the path looks fortunate only when judged by human values, and pointed to ``The Gift We Give To Tomorrow.''

In short: a summary that states its conclusion and places the argument in earlier posts; its examples show that removing one part of value can empty the future, but give no case of the lesser disturbances it also warns of.
\end{afterword}
